\begin{table}[t]
\caption{Robustez na MCC ($n=99.356$, $r=4.968$): remoção aleatória ($T=30$, média $\pm$ DP) e remoção dos $r$ vértices de maior grau. Sem remoção, C é o valor exato; após a remoção, C é estimada por BFS a partir de 100 (aleatória) e 2.000 origens (por grau).}
\label{tab:rob}
\footnotesize
\begin{tabular}{@{}lrrr@{}}
\toprule
Métrica & Sem remoção & Aleatória & Por grau \\
\midrule
A: maior componente & 99.356 & 92.804 $\pm$ 339 & 40.975 \\
A: $S$ & 1{,}0000 & 0{,}9832 $\pm$ 0{,}0036 & 0{,}4341 \\
B: componentes & 1 & 1.577{,}7 $\pm$ 337{,}8 & 52.451 \\
C: distância média & 3{,}725 & 3{,}724 $\pm$ 0{,}044 & 6{,}011 \\
C: maior distância & 9$^{\dagger}$ & 7{,}9 $\pm$ 0{,}4 & 17 \\
D: fração isolados & 0{,}0000 & 0{,}0166 $\pm$ 0{,}0036 & 0{,}5469 \\
\bottomrule
\multicolumn{4}{@{}l}{\footnotesize $^{\dagger}$ diâmetro exato; após a remoção, maior distância vista nas BFS.}
\end{tabular}
\end{table}
